% Part: proof-theory
% Chapter: natural-deduction
% Section: translation-G2i

\documentclass[../../../include/open-logic-section]{subfiles}

\begin{document}

\olfileid{pt}{nat}{tgi}

\olsection{الترجمة من \Log{G2i} إلى \Log{N2i}}

مقدّمات المتتاليات في~\Log{G2i} مجموعات متعددة~$\Gamma$ من
!!{formula}s، بينما مقدّمات المتتاليات في~\Log{N2i} مجموعات~$\Gamma'$ من
الصيغ الموسومة، بحيث لا يظهر كل وسم إلا مرة واحدة. لنقل إن مجموعة~$\Gamma'$
من~!!{formula}s الموسومة \emph{تقابل} المجموعة المتعددة~$\Gamma$ إذا كان، لكل
!!{formula}~$!A$، عدد~!!{element}s المجموعة~$\Gamma'$ التي صورتها
$x : !A$ أصغر من أو يساوي تعددية~$!A$ (أي عدد نسخ~$!A$) في~$\Gamma$.

\begin{prop}\ollabel{prop:G2i-to-N2i}
إذا كان $\Log{G2i} + \Cut \Proves \Gamma \Sequent \Delta$، فإن $\Log{N2i}
\Proves \Gamma' \Sequent !C$، حيث $!C \ident \lfalse$ إذا كانت $\Delta$
خالية، و$!C \ident !A$ إذا كانت $\Delta = \{!A\}$، وحيث تقابل $\Gamma'$
المجموعة~$\Gamma$.
\end{prop}

\begin{proof}
  نبيّن أنه إذا كان~$\pi$ !!a{proof} للمتتالية $\Gamma \Sequent \Delta$ في
  $\Log{G2i}+\Cut$، فثمة~$\Gamma'$ تقابل~$\Gamma$ وثمة
  !!a{proof}~$\delta$ للمتتالية~$\Gamma' \Sequent !C$ في~\Log{N2i}.
  نفعل ذلك بالاستقراء في $n = \pheight{\pi}$.

  إذا كان $n = 0$، فإن~$\pi$ مسلّمة. فإذا كانت المسلّمة $\lfalse \Sequent
  \quad$، كانت~$\delta$ هي المسلّمة $x:\lfalse \Sequent \lfalse$؛ أي إن
  $\Gamma' = \{x:\lfalse\}$ و$!C \ident \lfalse$. وإذا كانت~$\pi$ مسلّمة
  من الصورة $!A \Sequent !A$، كانت~$\delta$ هي $x: !A \Sequent !A$.
  
  وإذا كان $n>0$، انتهى~$\pi$ بقاعدة~$R$. يضمن فرض الاستقراء وجود
  !!{proof}s في~\Log{N2i} لمتتاليات تقابل مقدمات تلك القاعدة. وبإعادة
  الوسم عند الحاجة، نستطيع أن نفترض أن وسوم التفريغ والوسوم المقيّدة في هذه
  البراهين متميزة، وأنها متباينة بعضها عن بعض بين براهين المقدمات المختلفة
  أو المنسوخة، وأن الوسم~$x$ لا يظهر فوق استدلال يفرّغه إلا إذا كان ذلك
  الاستدلال يفرّغه بالفعل. ثم نبيّن كيف يمكن جمع هذه~!!{proof}s للحصول على
  براهين لمتتالية مناسبة~$\Gamma' \Sequent !C$. ونميّز الحالات بحسب~$R$.

  \begin{enumerate}
    \item إذا كانت~$R$ هي~\RightR{\Weakening}، انتهى~$\pi$ بما يأتي:
    \[
    \AxiomC{}
    \RightLabel{$\pi_1$}
    \Deduce$\Gamma \fCenter$
    \RightLabel{\RightR{\Weakening}}
    \UnaryInf$\Gamma \fCenter !A$
    \DisplayProof
    \]
    يعطي تطبيق فرض الاستقراء على~$\pi_1$ !!a{proof}
    للمتتالية~$\Gamma' \Sequent \lfalse$. ونطبّق~$\FalseInt$ لنحصل على
    !!a{proof} للمتتالية $\Gamma' \Sequent !A$.
    \item إذا كانت~$R$ هي~\LeftR{\Weakening}، انتهى~$\pi$ بما يأتي:
    \[
    \AxiomC{}
    \RightLabel{$\pi_1$}
    \Deduce$\Gamma \fCenter \Delta$
    \RightLabel{\LeftR{\Weakening}}
    \UnaryInf$!A, \Gamma \fCenter \Delta$
    \DisplayProof
    \]
    يعطي تطبيق فرض الاستقراء على~$\pi_1$ !!a{proof}
    للمتتالية~$\Gamma' \Sequent !C$، ونتخذه~$\delta$. لاحظ أنه إذا كانت
    $\Gamma'$ تقابل~$\Gamma$ فهي تقابل أيضًا~$!A,\Gamma$، لأن عدد الصيغ
    $x:!A$ في~$\Gamma'$ لا يزيد على عدد نسخ~$!A$ في~$\Gamma$، ومن ثم لا
    يزيد أيضًا على عدد نسخها في السياق~$!A,\Gamma$.
    \item إذا كانت~$R$ هي~\LeftR{\Contraction}، انتهى~$\pi$ بما يأتي:
    \[
    \AxiomC{}
    \RightLabel{$\pi_1$}
    \Deduce$!A, !A, \Gamma \fCenter \Delta$
    \RightLabel{\LeftR{\Contraction}}
    \UnaryInf$!A, \Gamma \fCenter \Delta$
    \DisplayProof
    \]
    يعطي تطبيق فرض الاستقراء على~$\pi_1$ !!a{proof}
    هو~$\delta_1$ للمتتالية~$\Pi, \Gamma' \Sequent !C$، حيث
    $\Pi \subseteq \{x : !A,y:!A\}$. إذا كانت
    $\Pi = \{x : !A,y:!A\}$، فاستبدل بكل~!!{formula}s الموسومة~$y:!A$ في
    $\delta_1$ الصيغة~$x:!A$. وبما أن~$x$ يظهر في سياق المتتالية النهائية
    لـ~$\delta_1$، يمكننا افتراض أن لا استدلال في~$\delta_1$ يفرّغ
    !!a{formula} من الصورة~$x:!B$؛ أي إن أي~$x:!B$ في يسار متتالية ضمن
    $\delta_1$ ليس فعالًا. ولذلك تظل النتيجة~!!a{proof} في~\Log{N2i}.
    \item إذا كانت~$R$ هي~$\RightR{\lif}$، انتهى~$\pi$ بما يأتي:
    \[
    \AxiomC{}
    \RightLabel{$\pi_1$}
    \Deduce$!A, \Gamma \fCenter !B$
    \RightLabel{\RightR{\lif}}
    \UnaryInf$\Gamma \fCenter !A \lif !B$
    \DisplayProof
    \]
    يوجد، بفرض الاستقراء، !!a{proof}~$\delta_1$ للمتتالية
    $x:!A,\Gamma' \Sequent !B$ أو للمتتالية~$\Gamma' \Sequent !B$، حيث
    تقابل~$\Gamma'$ المجموعة~$\Gamma$. نحصل على~$\delta$ بإتباع
    $\delta_1$ بتطبيق~\Intro{\lif} (ويجوز أن يكون الفرض الموسوم غير
    مستعمل).
    \item إذا كانت~$R$ هي~\LeftR{\lif}، انتهى~$\pi$ بما يأتي:
    \[
    \AxiomC{}
    \RightLabel{$\pi_1$}
    \Deduce$\Gamma_1 \fCenter !A$
    \AxiomC{}
    \RightLabel{$\pi_2$}
    \Deduce$!B, \Gamma_2 \fCenter \Delta$
    \RightLabel{\LeftR{\lif}}
    \BinaryInf$!A \lif !B, \Gamma_1, \Gamma_2 \fCenter \Delta$
    \DisplayProof
    \]
    يعطي فرض الاستقراء !!{proof}s: برهانًا~$\delta_1$ للمتتالية
    $\Gamma_1' \Sequent !A$، وبرهانًا~$\delta_2$ للمتتالية
    $x:!B,\Gamma_2' \Sequent !C$ أو للمتتالية~$\Gamma_2' \Sequent !C$.
    في الحالة الثانية يصلح~$\delta_2$ أن يكون~$\delta$. وفي الحالة الأولى
    نعيد تسمية جميع الصيغ ووسوم التفريغ والوسوم المقيّدة، تجنبًا للأسر، بحيث
    لا يظهر أي وسم في كل من~$\delta_1$ و$\delta_2$. عندئذ يقابل السياق
    $\Gamma_1',\Gamma_2'$ السياق~$\Gamma_1,\Gamma_2$. ولأن~$x:!B$ يظهر في
    المتتالية النهائية لـ~$\delta_2$، فلا يمكن أن يكون~$x:!B$ فعالًا في أي
    استدلال وسم تفريغه~$x$ داخل~$\delta_2$. استبدل بكل مسلّمة
    $x:!B \Sequent !B$ في~$\delta_2$ الـ~!!{proof}~$\delta_3$ الآتي:
    \[
    \Axiom$x: !A \lif !B \fCenter !A \lif !B$
    \AxiomC{}
    \RightLabel{$\delta_1$}
    \Deduce$\Gamma_1' \fCenter !A$
    \RightLabel{\Elim{\lif}}
    \BinaryInf$x: !A \lif !B, \Gamma_1' \fCenter !B$
    \DisplayProof
    \]
    واستبدل بكل ظهور لـ~$x:!B$ في~$\delta_2$ الصيغة~$x:!A \lif !B$.
    والنتيجة~$\Subst{\delta_2}{\delta_3}{x:!B}$ هي~!!a{proof} للمتتالية
    $x:!A \lif !B,\Gamma_1',\Gamma_2' \Sequent !C$.
    \item إذا كانت~$R$ هي~\RightR{\land}، انتهى~$\pi$ بما يأتي:
    \[
    \AxiomC{}
    \RightLabel{$\pi_1$}
    \Deduce$\Gamma_1 \fCenter !A$
    \AxiomC{}
    \RightLabel{$\pi_2$}
    \Deduce$\Gamma_2 \fCenter !B$
    \RightLabel{\RightR{\land}}
    \BinaryInf$\Gamma_1, \Gamma_2 \fCenter !A \land !B$
    \DisplayProof
    \]
    يعطي فرض الاستقراء برهانًا~$\delta_1$ للمتتالية
    $\Gamma_1' \Sequent !A$ وبرهانًا~$\delta_2$ للمتتالية
    $\Gamma_2' \Sequent !B$. وبإعادة تسمية جميع الصيغ ووسوم التفريغ والوسوم
    المقيّدة في~$\delta_2$ تجنبًا للأسر، يمكننا ضمان ألا يظهر أي وسم في كل من
    $\delta_1$ و$\delta_2$. وعندئذ يقابل السياق~$\Gamma_1',\Gamma_2'$
    السياق~$\Gamma_1,\Gamma_2$. ونحصل على~!!a{proof}~$\delta$ بتطبيق
    $\Intro{\land}$ على المتتاليتين النهائيتين لـ~$\delta_1$ و$\delta_2$.
    \item إذا كانت~$R$ هي~$\LeftR{\land}$، وانتهى~$\pi$، مثلًا، بما يأتي:
    \[
    \AxiomC{}
    \RightLabel{$\pi_1$}
    \Deduce$!A, \Gamma \fCenter \Delta$
    \RightLabel{\LeftR{\land}}
    \UnaryInf$!A \land !B, \Gamma \fCenter \Delta$
    \DisplayProof
    \]
    فبفرض الاستقراء يوجد~!!a{proof}~$\delta_1$ للمتتالية
    $x:!A,\Gamma' \Sequent !C$ أو للمتتالية~$\Gamma' \Sequent !C$، حيث
    تقابل~$\Gamma'$ المجموعة~$\Gamma$. في الحالة الثانية نتخذ~$\delta_1$
    ليكون~$\delta$. وفي الحالة الأولى، لاحظ أن ظهور~$x:!A$ في المتتالية
    النهائية يعني أنه ليس فعالًا في أي استدلال وسم تفريغه~$x$ داخل
    $\delta_1$. استبدل بكل مسلّمة $x:!A \Sequent !A$ الـ~!!{proof}~$\delta_2$
    الآتي:
    \[
    \Axiom$x: !A \land !B \fCenter !A \land !B$
    \RightLabel{\Elim{\land}[1]}
    \UnaryInf$x: !A \land !B \fCenter !A$
    \DisplayProof
    \]
    واستبدل بكل ظهور لـ~$x:!A$ الصيغة~$x:!A \land !B$؛ أي انظر في
    $\Subst{\delta_1}{\delta_2}{x:!A}$. وهذا~!!a{proof} للمتتالية
    $x:!A \land !B,\Gamma' \Sequent !C$.
    \item إذا كانت~$R$ هي~\Cut، انتهى~$\pi$ بما يأتي:
    \[
    \AxiomC{}
    \RightLabel{$\pi_1$}
    \Deduce$\Gamma_1 \fCenter !A$
    \AxiomC{}
    \RightLabel{$\pi_2$}
    \Deduce$!A, \Gamma_2 \fCenter \Delta$
    \RightLabel{\Cut}
    \BinaryInf$\Gamma_1, \Gamma_2 \fCenter \Delta$
    \DisplayProof
    \]
    يعطي فرض الاستقراء برهانًا~$\delta_1$ للمتتالية
    $\Gamma_1' \Sequent !A$ وبرهانًا~$\delta_2$ للمتتالية
    $x:!A,\Gamma_2' \Sequent !C$ أو للمتتالية~$\Gamma_2' \Sequent !C$.
    نعيد وسم نسختي البرهان تجنبًا للأسر بحيث تكون وسومهما متباينة. وفي
    الحالة الثانية نتخذ~$\delta=\delta_2$. وإلا استبدل بكل مسلّمة
    $x:!A \Sequent !A$ في~$\delta_2$ البرهان~$\delta_1$ لنحصل على~$\delta$
    بوصفه~!!{proof}~$\Subst{\delta_2}{\delta_1}{x:!A}$ للمتتالية
    $\Gamma_1',\Gamma_2' \Sequent !C$.
    \item تُعالج الحالات المنطقية المتبقية على نحو مماثل.
  \end{enumerate}
\end{proof}

\end{document}
